\documentclass[11pt]{article}

% --- Layout: kompakt, für ca. 1 Seite Stichpunkte ---
\usepackage[a4paper,margin=2cm]{geometry}
\usepackage[ngerman]{babel}
\usepackage[utf8]{inputenc}
\usepackage[T1]{fontenc}
\usepackage{enumitem}
\usepackage{titlesec}
\usepackage{hyperref}
\usepackage{booktabs}
\usepackage{csquotes}
\hypersetup{colorlinks=true, linkcolor=black, urlcolor=blue}

\setlist{nosep, leftmargin=1.4em, itemsep=1pt, topsep=2pt}
\titlespacing*{\section}{0pt}{8pt}{3pt}
\titlespacing*{\subsection}{0pt}{5pt}{2pt}
\titleformat{\section}{\normalfont\bfseries\large}{\thesection}{0.6em}{}
\titleformat{\subsection}{\normalfont\bfseries}{\thesubsection}{0.6em}{}

\pagenumbering{gobble} 

\begin{document}

\begin{center}
    {\Large\bfseries Notizen --- Vorab-Aufgabe Datenkonsolidierung}\\[2pt]
    \small [Kristian] \hfill [16.08.2026]
\end{center}
\vspace{4pt}
% =========================================================
\section{Prioritäten}
\begin{enumerate}
    \item Zuverlässiges, robustes und reproduzierbares Vorgehen zur Extraktion und Migration der notwendigen Informationen.
    \item Inkrementelles Erweitern des Outputs der Felder um sich vorläufigem Ziel-Schema anzunähern. Grundlage hierfür eine zukunftsfähige und wartbare Lösung für kommerziellen Gebrauch.
    \item Erarbeiten besonderer Herausforderungen der Datenkonsolidierung (z.B. fehlende Informationen, Inkonsistenzen, Abweichungen zwischen Standorten).
    \item Transparente Dokumentation der Entscheidungen und Vorgehensweise, um die Nachvollziehbarkeit zu gewährleisten. Auflistung weiterer Potenziale und aktueller Leistungen.
\end{enumerate}

% =========================================================
\section{Vorgehen: Ziel 1}
\subsection{Exploration}
\begin{itemize}
    \item (I) Prüfen der Intra-Standort Variabilität der Eingaben (i.e. wie hoch ist die Eingabekonsistenz innerhalb der Standorte)
    \item (I) Prüfen der Inter-Standort Variabilität der Eingaben (i.e. wie hoch ist die Eingabekonsistenz zwischen den Standorten)
    \item Ziel: Gibt es Gemeinsamkeiten zwischen den Standorten die die eine eineitliche Extraktion von Informationen zulassen (\"required\")?
    \item Ziel: Welche Abweichungen erschweren die Extraktion von Informationen?
    \item Ergebnis: Folgende Tabelle zeigt welche Informationen einheitlich aus allen Dateien zuverlässig ausgelesen werden können:
\end{itemize}
\begin{table}[h]
    \centering
    \begin{tabular}{@{} p{0.32\linewidth} p{0.58\linewidth} @{}}
        \toprule
        \textbf{.html Quelle} & \textbf{Info} \\
        \midrule
        Titel (<h1> Tag) & Zimmer, Vermarktungsart, Objekttyp \\
        <div class=\enquote{expose} data-objekt-nr=\enquote{xxxx}> & Objekt-ID \\
        Objekt-Pfad & Quelle \\
        <p class=\enquote{lage}>\enquote{xxxx}</p> & Adresse\\
        <table class=\enquote{eckdaten}> & Wohnflaeche\\
        \bottomrule
    \end{tabular}
\end{table}
\subsection{Umsetzung der Erkenntisse}
\begin{itemize}
    \item Regex und Wortsuche aus auf <h1> als Primärquelle mit Fallback auf Eckdaten Tabelle. Warnungen und Errors werden geloggt und gesammelt mit einer Fehlerbeschreibung.
    \item Aufgrund von nicht eideutig lösbaren Inkonsistenzen erwarten die fehlerhaften Outputs einen HITL Ansatz.
\end{itemize}
% =========================================================


% =========================================================


% =========================================================


% =========================================================


\end{document}